\section{The Case for Self-Regulation}
\label{sec:vision}

The central claim is not that vehicles should form a fully cooperative swarm. It is that each vehicle should remain the sole decision maker for its own motion, with network-level behavior emerging from repeated local damping, not from an agreed-upon plan. This distinction matters for deployment. Command-style coordination creates hard questions about authority, liability, and failure: who tells whom to slow down, and what if the message is late? Self-regulation avoids that contract, asking nothing of neighboring vehicles beyond their physical presence.

\noindent\textbf{Why local damping suffices.}
Highway regulation does not require any vehicle to know the network state, only early evidence that local flow is nearing a bad operating region: density rising near a merge, queues forming downstream, speed variance growing before a shockwave. The control objective is physical and local: keep density below the fundamental diagram's critical point, then hold speed high and variance low. A vehicle that damps its own speed variance when it sees rising local pressure removes energy from an incipient wave before it steepens, and when enough vehicles do this independently the stream is held near critical density with no shared plan. Even a single uncoordinated vehicle, with no communication at all, already damps stop-and-go waves on ring tracks and open roads~\cite{stern2018dissipation,lee2024traffic}, so decentralized local action is not the speculative part. The property to protect as actuators are added is string stability: disturbances must not amplify as they propagate upstream through a chain of followers~\cite{swaroop1996string}, which gives the safety layer and evaluation a concrete target. The analogy to backpressure scheduling, which stabilizes queues from local pressure differences alone, is direct~\cite{tassiulas1990stability}; but packets do not brake hard when startled, which is why Section~\ref{sec:architecture} pairs the controller with a safety filter.

\noindent\textbf{Communication model.}
Better-than-myopic observability is where communication earns its place, and we keep its role thin. Participating vehicles exchange compact, anonymous, aggregate summaries: nearby counts and density for participants and ordinary traffic, mean speed, lane distribution, a local queue estimate, downstream congestion, and merge pressure. They never exchange raw video, stable identities, acceleration or braking commands, or joint lane plans. Three receiver-side semantics matter more than any message format: expired peers are dropped rather than extrapolated, unheard peers map to explicit neutral values rather than zeros with hidden meaning, and the aggregate is advisory evidence, never a command. Our prototype realizes the substrate as 2\,Hz broadcasts with a 150\,m radius and a 2\,s time-to-live; a summary is tens of bytes, so the footprint is kilobits per second and does not grow with scene complexity. Any deployed radio can carry it, whether DSRC, C-V2X, or in the limit a cellular relay~\cite{sae2020j2735,etsi2019cam,kenney2011dedicated}. If no peer is heard, a vehicle degrades to a local-only controller. Because the consumed unit is a neighborhood aggregate, no camera frames, identifiers, or trajectories need be published for neighbors to benefit; realizing that guarantee end to end is an open problem (Section~\ref{sec:discussion}).

\noindent\textbf{Minimal fleets at bottlenecks.}
The usual objection to vehicle-based control is adoption: no consumer rollout reaches meaningful market share quickly. But market share is the wrong unit. Congestion is not uniform; it is manufactured at bottlenecks, at merges, lane drops, and tunnels, and the network's worst behavior is set by the inflow those few places receive. The relevant metric is presence in a specific traffic stream at a specific place, which a single fleet operator controls today. Consider the Holland Tunnel, which carries Interstate~78 under the Hudson through four lanes and roughly $10^5$ vehicles per day~\cite{panynj2024traffic}. Every eastbound trip is funneled through approach corridors totaling about 25 feeder lanes, and since adding lanes under the river is impossible, the only lever is the inflow. Ramp metering is exactly this idea, but it requires owning the ramp; a self-regulating car can be introduced one vehicle at a time. The sparse-actuator result suggests one controlled vehicle per roughly twenty damps waves in its neighborhood~\cite{stern2018dissipation}, so one self-regulating vehicle in the platoon window of each feeder lane, about 25 present at a time, is the scale at which the tunnel's inflow becomes controllable. A single rideshare, parcel, or telematics-insured fleet already has that ambient presence at peak, and the operator capturing the benefit, predictable trip times through the bottleneck, is the one paying the actuation cost. The predecessor points the same way: it modulated speed only at bottleneck segments during peak inflow, and its mixed-traffic ablation improves steadily with the compliant fraction rather than requiring full adoption~\cite{self_regulating_cars}. The aggregate signal cooperates with this deployment shape: the number of peers a vehicle hears scales with the very density it watches. Near critical density, a 150\,m window on a three-lane highway holds roughly 22 vehicles, about one cooperating peer at 5\% penetration, and several near jam density. The honest floor: near free flow at very low penetration a vehicle often hears nothing and falls back to local-only control. Where on this penetration--density curve the aggregate begins to change outcomes, and whether bottleneck seeding beats uniform adoption at equal fleet size, are central empirical questions (Sections~\ref{sec:eval} and~\ref{sec:discussion}).
